\section{Обзор существующих решений}
\label{sec:Chapter2} \index{Chapter2}

%Здесь надо рассмотреть все существующие решения поставленной задачи, но не
%просто пересказать, в чем там дело, а оценить степень их соответствия тем
%ограничениям, которые были сформулированы в постановке задачи.

Есть несколько широко известных $\mathsf{WCH}$ проблем на решётках: $\mathsf{SVP}$ ({\it Shortest Vector Problem}), $\mathsf{GapSVP}$, $\mathsf{SIVP}$ ({\it Shortest independent vectors problem}) и т.п. В $1996$ году Ажтай свёл часть из них к $\mathsf{ACH}$ проблеме $\mathsf{SIS}$ ({\it Short Integer Solution})~\cite{ACH_Ajtai}. Позже Регев проделал аналогичную вещь с $\mathsf{LWE}$.
